% =====================================================================
%  第 4 节：问题为何未解决
% =====================================================================
\section{问题为何未解决}
\label{sec:openness}

问题 \ref{prob:main} 不是一个已被解决的问题，也不是一个“工程实现问题”。本节说明它难在哪里——分为\textbf{数学层面的困难}（\S\ref{sec:hard-math}）、\textbf{统计层面的困难}（\S\ref{sec:hard-stat}）、以及\textbf{方法论层面的困难}（\S\ref{sec:hard-method}）。把困难讲清楚，也是说明本课题为何“允许不彻底解决”的前提。

\subsection{数学层面的困难}
\label{sec:hard-math}

\subsubsection{候选空间的开放性与可读性的矛盾}

若不加限制，猜想空间是无穷的：可以引入任意复合不变量（如 $\chi^2/\omega$、$\log\alpha$、$\Delta^{\alpha}$），可以加入任意取整与条件分支。理论上，越大的空间越可能包含“有趣的”猜想，但同时也使候选的可读性与可证性急剧下降。

这里存在一个\textbf{实质性的张力}：

\begin{itemize}
  \item 限制在\textbf{线性、小整数系数、稀疏}的形式内，候选可读、可用 LP 求解、可用初等方法验证——但已知的大量有趣猜想（Reed 猜想的取整、Nordhaus–Gaddum 型的两图关系）落在这个空间之外。
  \item 放开到\textbf{非线性、含取整}的形式，表达力足够——但“求解系数”不再是 LP，而变成混合整数非线性规划，且“这个猜想是否已知”的判断也变得困难。
\end{itemize}

本课题采取折中：\textbf{以线性形式为主干（可自动化），用 LLM 提议的非线性形式作为补充（半自动化）}。非线性候选不进入 LP，而由 LLM 给出系数后直接检验。这使系统能触及 Reed 猜想这类形式，但也意味着这部分候选的质量控制更依赖 LLM，需要单独评估。

\subsubsection{“已知”的判定问题}

一条候选是否“新”，没有算法化的判据。数学文献中的等价变形（移项、乘以正数、代入已知不等式）可以产生无穷多种表述同一事实的方式。判断“$\chi+\alpha\le n+1$”与“$\alpha\le n-\chi+1$”是同一个定理，需要语义层面的理解。

这构成一个\textbf{无法完全自动化}的环节。本课题的做法是：维护一个已知结论的清单（数十条经典不等式），用 LLM 做语义匹配，并\textbf{明确标注}“新颖性判断是启发式的、可能出错”。在最终报告中，所有“候选新猜想”都会附上人工核对的结果。

\subsubsection{反例的存在性不可判定}

对一般的 $\bm{c}$，判定“是否存在图使 $\bm{c}^\top\bm{f}>0$”在理论上是困难的：$\G$ 无穷，且 $\alpha,\omega,\chi$ 的计算本身是 NP-难的。即使限定 $n\le N$，枚举 $\G_N$ 的代价也随 $N$ 指数增长（$|\G_{10}|\approx1.2\times10^7$，$|\G_{11}|\approx10^9$）。

因此在实践中，反驳永远是\textbf{不完整的}：我们只能在一个可及的范围内搜索。这与“证明”有本质区别——证明是对全体的断言，搜索只是对部分的检查。

\subsection{统计层面的困难}
\label{sec:hard-stat}

\subsubsection{零分布未知}

\S\ref{sec:f-screen} 提出的置换检验方案，其有效性依赖于一个未经验证的假设：置换后得到的 $\tilde{T}_j$ 分布确实近似于 $H_0$ 下的分布。问题在于：

\begin{itemize}
  \item 不变量之间存在\textbf{确定性约束}（$\omega\le\chi$、$\chi\le\Delta+1$、$\alpha\omega\ge n$ 等）。随机置换会破坏这些约束，产生的“伪图”不变量组合可能不对应任何真实图，从而使零分布偏离真实情形。
  \item 候选之间\textbf{强相关}。同一个不变量出现在多条候选中，一条被推翻往往意味着相关的几条也被推翻。BH 在任意相依下不保证 FDR 控制。
\end{itemize}

这些问题的严重程度\textbf{只能通过实验评估}。我们计划用“已知定理 + 已知反例”构成的有标签集合校准：若某策略在该集合上的 FDR 显著超过名义水平 $q$，则说明零分布构造有误，需要修正（例如改用 \S\ref{sec:f-refute} 中记录的“搜索充分性”作为协变量做条件检验）。

\subsubsection{预算分配引入的选择偏差}

汤普森采样把预算集中到“看起来可疑”的候选上，这本身是合理的——但它使不同候选的“未被推翻”具有\textbf{不同的证据强度}。被少搜的候选，其 $T_j$ 更可能是“没搜到”而非“真的没有”。

这是一个典型的\textbf{选择性推断}（selective inference）问题：数据依赖的决策（把预算给谁）破坏了后续检验的分布性质。文献中的处理方式包括样本分割（用一部分预算决定分配、另一部分做检验）与条件推断（在给定分配的条件下做检验）。本课题倾向于\textbf{样本分割}，因为它简单且不需要推导复杂的条件分布，代价是浪费一部分预算。

\subsubsection{外推的系统性偏移}

$n\le 9$ 的图与 $n=12$ 的图不只是“大小不同”，而是\textbf{结构上不同}。例如：

\begin{itemize}
  \item 需要较多顶点才能实现的构造（如某些 Mycielski 迭代、Kneser 图 $K(5,2)$ 需 10 个顶点）在 $n\le 9$ 中完全缺席；
  \item 小图上“几乎所有图都是连通的、度分布集中”，而大图上存在更多极端结构；
  \item 某些不等式在小图上因“没有足够的空间构造反例”而成立，在大图上则失效。
\end{itemize}

因此从发现集到测试集的外推，其偏移是\textbf{系统性}的，不能靠增大样本量消除。这构成本课题最根本的方法论限制，也是我们在结论中必须明确声明的。

\subsection{方法论层面的困难}
\label{sec:hard-method}

\subsubsection{“成功”的定义本身是模糊的}

一个自动猜想系统怎样算成功？可能的判据包括：

\begin{itemize}
  \item 提出的猜想中有多少后来被证明（需要数年时间验证，无法在课程周期内评估）；
  \item 提出的猜想中有多少“有意思”（主观）；
  \item 系统能否重现已知定理（可评估，但重现已知不算发现）；
  \item 系统能否产出可被人类数学家在一小时内判断真伪的猜想（可评估，且实用）。
\end{itemize}

本课题采用最后一条作为主要判据，辅以第三条作为正确性检验。这个选择本身是一个\textbf{方法论主张}：自动猜想系统的价值在于\textbf{降低人类筛选成本}，而不在于替代人类做证明。

\subsubsection{LLM 的不稳定性}

LLM 的输出具有随机性：同一个提示在不同时刻可能给出不同的形式提议。这带来两个问题：

\begin{itemize}
  \item \textbf{可复现性}：整个循环的结果依赖于 LLM 的采样。我们通过记录每次调用的完整输入输出（写入轨迹文件）、并固定随机种子与模型版本来缓解，但无法完全消除。
  \item \textbf{评估的混淆}：若最终结果不理想，难以判断是“方法不行”还是“这次 LLM 恰好表现差”。需要通过多次独立运行（不同种子）报告结果的分布，而非单次结果。
\end{itemize}

\subsubsection{“智能体”与“脚本”的界限}

这是本课题必须自我审视的一点。一个自然的质疑是：既然猜想生成可以用 LP、反驳可以用搜索、筛选可以用 BH，那么 LLM 在其中是否只是装饰？若去掉 LLM 系统仍能运行，是否说明它不构成“智能体”？

我们的回应是：LLM 承担的是\textbf{在开放空间中的定向搜索}——提议哪些形式值得试、根据失败原因修正形式、判断候选是否已知。这些环节的规则无法预先枚举（否则就退化为穷举），需要语义理解。但这一主张需要\textbf{实验支持}：我们计划做消融实验，比较“有 LLM 提议”与“随机/穷举提议”在相同预算下的产出质量。若两者无显著差异，则应诚实报告“LLM 在此任务中贡献有限”——这同样是一个有价值的结论。

\begin{notebox}[关于本课题的定位]
上述困难说明：本课题\textbf{不可能}在课程周期内被彻底解决。它的合理目标是——搭建一个可运行的 CRP 循环，在一个受限但明确的设定下（线性形式、$n\le 12$、已知不变量集合）产出可审计的结果，并对“统计筛选是否有效”给出初步的实证回答。这符合课程“允许不彻底解决、但须体现运行与迭代”的要求。
\end{notebox}
